\documentclass[12pt]{article}
%---DOCUMENT MARGINS---
\usepackage{geometry} % Required for adjusting page dimensions and margins
\geometry{
	paper=a4paper, % Paper size, change to letterpaper for US letter size
	top=4cm, % Top margin
	bottom=2cm, % Bottom margin
	left=2cm, % Left margin
	right=2cm, % Right margin
	headheight=3cm, % Header height
	%footskip=1.5cm, % Space from the bottom margin to the baseline of the footer
	headsep=0.5cm, % Space from the top margin to the baseline of the header
	%showframe, % Uncomment to show how the type block is set on the page
}
\usepackage{cmap}

\usepackage[T1,T2A]{fontenc}
\usepackage{fontspec}
\setmainfont[
	BoldFont={* Bold},
	ItalicFont={* Italic},
	BoldItalicFont={* BoldItalic}
]{DejaVu Serif Condensed}
\setsansfont{DejaVu Sans}
\setmonofont{Dejavu Sans Mono}
\usepackage[math-style=TeX]{unicode-math}
\setmathfont{Latin Modern Math}

\usepackage[nobottomtitles*]{titlesec}
\titleformat
{\section} % command
{\fontsize{14}{14}\sffamily\bfseries} % format
{} % label
{0pt} % sep
{} % before-code
\titlespacing{\section}{0pt}{0.75em}{0.25em}
\newcommand{\tl}{}
\newcommand{\ml}{}
\newcommand{\problem}[1]{%
	\section[#1]{#1 \fontsize{14}{14}\selectfont{\hfill{\normalfont{
				\emoji{hourglass-not-done}\tl \space\space \emoji{floppy-disk} \ml}}}}
}
\AddToHook{cmd/section/before}{\clearpage}

\usepackage[dvipsnames]{xcolor}
\titleformat
{\subsection} % command
{\fontsize{14}{14}\sffamily\bfseries} % format
{\includesvg[width=0.035\textwidth, inkscapelatex=false]{lion}\space} % label
{0pt} % sep
{} % before-code
\titlespacing{\subsection}{0pt}{0.75em}{0.25em}

\setlength{\parskip}{0.5em}
\setlength{\parindent}{0pt}
\renewcommand{\baselinestretch}{1.2}
\sloppy

\usepackage{listings}
\definecolor{lstbg}{RGB}{250,250,252}
\definecolor{lstframe}{RGB}{220,220,225}
\lstdefinestyle{mystyle}{
	backgroundcolor=\color{lstbg},
	frame=single,
	rulecolor=\color{lstframe},
	basicstyle=\ttfamily\small,
	keywordstyle=\bfseries,
	breaklines=true,
	aboveskip=0.5\baselineskip,
	belowskip=0\baselineskip  
}
\lstset{style=mystyle, language=C++}

\usepackage{fancyhdr}
\pagestyle{fancy}
\setlength{\headwidth}{\textwidth}
\newcommand{\round}{}
\newcommand{\problemName}{}
\newcommand{\lang}{English}
\fancyhead[L]{
	\begin{minipage}{0.4\textwidth}
		\bf{
			EJOI 2025 \round\\
			\problemName\\
			\emoji{globe-with-meridians} \lang
		}
	\end{minipage}
	\vspace{0.5cm}
}
\fancyhead[R]{%
	\begin{minipage}{0.6\textwidth}
		\includesvg[width=\textwidth, inkscapelatex=false]{logo}
	\end{minipage}
	\vspace{0.5cm}
}
\usepackage{lastpage}
\fancyfoot[C]{\problemName\space(\lang) \thepage\ / \pageref*{LastPage}}

\raggedbottom

\usepackage{amsmath}
\usepackage{stmaryrd}

\usepackage{graphicx}
\graphicspath{{./}}
\usepackage[export]{adjustbox}
\usepackage{wrapfig}
\makeatletter
\patchcmd\WF@putfigmaybe{\lower\intextsep}{}{}{\fail}
\AddToHook{env/wrapfigure/begin}{\setlength{\intextsep}{0pt}}
\makeatother
\usepackage[inkscapearea=page, inkscapepath=./svg-inkscape]{svg}
\svgpath{{./}}

\usepackage{makecell}
\usepackage{tabularray}
\AtBeginEnvironment{table}{\vspace{-0.2cm}}
\AtEndEnvironment{table}{\vspace{-0.2cm}}
\usepackage{float}
\usepackage{placeins}
\usepackage{caption}
\captionsetup[table]{
	skip=0.25em,
	singlelinecheck=false, justification=justified
}

\usepackage{enumitem}
\setlist{itemsep=-0.4em, leftmargin=22pt, topsep=-\parskip}
\AfterEndEnvironment{itemize}{\vspace{\parskip}}
\newcommand{\tabitem}{\indent~~\llap{\textbullet}~~}

\usepackage{hyperref}
\hypersetup{
	colorlinks=true,
	citecolor=blue,
	linkcolor=blue,
	urlcolor=cyan,
}

\usepackage{emoji}

\usepackage[english]{babel}

\begin{document}

\renewcommand{\round}{Ден 1}
\renewcommand{\problemName}{Задача Затвор}
\renewcommand{\lang}{Македонски}

\renewcommand{\tl}{$2$ sec.}
\renewcommand{\ml}{$1024$ MB}
\problem{\problemName}

Иван и Петар се незаконски затворени во затвор. Сега тие мора да го испланираат нивното бекство. За да го направат тоа, тие мора да комуницираат што е можно поефикасно (всушност, Иван треба да му испраќа дневни информации на Петар). Но, тие неможе да се на иста локација во исто време и може само да комуницираат преку пораки напишани на паломи. Секој ден Иван сака да му испрати нова информација на Петар -- еден број помеѓу $0$ и $N-1$. Секој ден, за време на ручекот, Иван собира три паломи и запишува еден број помеѓу $0$ и $M-1$ на секоја палома (може истиот број да се запише на повеќе паломи) и ги остава на својот стол. Потоа, нивниот непријател, Бидик, секој ден краде по една палома (остануваат само две паломи на столот на Иван), и може да го смени редоследот на паломите (од останатите две). На крај, Петар ги наоѓа преостанатите две паломи и ги чита броевите запишани на нив. Тој мора точно да го одреди оригиналниот број кој Иван сакал да му го пренесе. Паломите имаат ограничен простор, па $M$ е фиксирана константа. Но, целта на Иван и Петар е да го максимизираат множеството на броеви кои ќе ги користат во размената на информации, па тие се слободни да ја одберат вредноста на $N$. Помогнете им на Иван и Петар така што ќе имплементирате стратегија за секој од нив, така што ќе ја максимизирате вредноста на $N$.


\subsection{Имплементациски детали}
Бидејќи ова е комуникациски проблем, вашата програма ќе се изврши два пати (еднаш за Иван и еднаш за Петар), тие неможат да памтат податоци или да комуницираат на било кој друг начин, освен начинот објаснат подолу. Вие треба да имплементирате три функции:

\begin{lstlisting}
 int setup(int M);
\end{lstlisting}


Оваа функција ќе биде повикана еднаш за време на извршувањето на Иван и еднаш за време на извршувањето на Петар. На влез ви е дадена вредноста на $M$ и мора да ја врати вредноста на $N$, двете извршувања на \texttt{setup} мора да ја вратат истата вредност на $N$.

\begin{lstlisting}
 std::vector<int> encode(int A);
\end{lstlisting}

Тука се имплементира стратегијата на Иван. Оваа функција ќе биде повикана со еден број $A$ ($0 \leq A < N$) и мора да врати три цели броеви $W_1, W_2, W_3$ ($0 \leq W_i < M$). Оваа функција ќе биде повикана вкупно $T$ пати -- еднаш за секој ден (вредности на $A$ може да се повторуваат помеѓу денови).


                                                        
\begin{lstlisting}
 int decode(int X, int Y);
\end{lstlisting}

Тука се имплементира стратегијата на Петар. Оваа функција ќе биде повикана со два од три броеви кои се добиени од \texttt{encode} во некој редослед. Мора да ја врати истата вредност на $A$ која била пратена до \texttt{encode}. Оваа функција ќе биде повикана вкупно $T$ пати -- според истите повици до функцијата \texttt{encode}; повиците ќе бидат во истиот редослед. Сите повици до \texttt{encode} ќе се случат пред сите повици до \texttt{decode}.

\subsection{Ограничувања}

\begin{itemize}
\item $M \leq 4300$
\item $T = 5000$
\end{itemize}

\subsection{Оценување}

За одредена подзадача, бројот на поени кои го добивате зависи од најмалата вредност на $N$ која била вратена од \texttt{setup} за било кој тест пример во подзадачата. Исто така зависи од $N^*$, која е таргет вердноста на $N$ за која морате точно да ја решите задачата за да ги добиете сите поени за подзадачата:

\begin{itemize}
\item Ако вашето решение има грешка на било кој тест пример, тогаш $S = 0$.
\item Ако $N \geq N^*$, тогаш $S = 1.0$.
\item Ако $N < N^*$, тогаш $S = \max\left(0.35 \max\left(\frac{\log(N) - 0.985 \log(M)}{\log(N^*) - 0.985 \log(M)}, 0.0\right)^{0.3} + 0.65 \left(\frac{N}{N^*}\right)^{2.4}, 0.01\right)$.
\end{itemize}

\subsection{Подзадачи}

\begin{table}[H]
\begin{tblr}{|Q[c,m]|Q[c,m]|Q[c,m]|Q[c,m]|}
\hline
\textbf{Подзадача} & \textbf{Поени} & $M$ & $N^*$ \\
\hline
$1$ & $10$ & $700$ & $82017$ \\
\hline
$2$ & $10$ & $1100$ & $202217$ \\
\hline
$3$ & $10$ & $1500$ & $375751$ \\
\hline
$4$ & $10$ & $1900$ & $602617$ \\
\hline
$5$ & $10$ & $2300$ & $882817$ \\
\hline
$6$ & $10$ & $2700$ & $1216351$ \\
\hline
$7$ & $10$ & $3100$ & $1603217$ \\
\hline
$8$ & $10$ & $3500$ & $2043417$ \\
\hline
$9$ & $10$ & $3900$ & $2536951$ \\
\hline
$10$ & $10$ & $4300$ & $3083817$ \\
\hline
\end{tblr}
\end{table}
\FloatBarrier

\subsection{Пример}

Да го разгледаме следниот пример со $T = 5$. Тука има стратегија така што Иван праќа три исти броеви за да ја енкодира вредноста $0$ и три различни вредности за да ја енкодира вредноста $1$. Забележете дека Петар може да ги декодира оригиналните броеви од било кои два од трите броеви кои Иван ги запишал.

\begin{table}[H]
\begin{tblr}{|Q[c,m]|Q[c,m]|Q[c,m]|}
\hline
\textbf{\makecell{Execution}} & \textbf{\makecell{Function call}} & \textbf{Return value} \\
\hline
Alice & \texttt{setup(10)} & \texttt{2} \\
\hline
Bob & \texttt{setup(10)} & \texttt{2} \\
\hline
Alice & \texttt{encode(0)} & \texttt{\{5, 5, 5\}} \\
\hline
Alice & \texttt{encode(1)} & \texttt{\{8, 3, 7\}} \\
\hline
Alice & \texttt{encode(1)} & \texttt{\{0, 3, 1\}} \\
\hline
Alice & \texttt{encode(0)} & \texttt{\{7, 7, 7\}} \\
\hline
Alice & \texttt{encode(1)} & \texttt{\{6, 2, 0\}} \\
\hline
Bob & \texttt{decode(5, 5)} & \texttt{0} \\
\hline
Bob & \texttt{decode(8, 7)} & \texttt{1} \\
\hline
Bob & \texttt{decode(3, 0)} & \texttt{1} \\
\hline
Bob & \texttt{decode(7, 7)} & \texttt{0} \\
\hline
Bob & \texttt{decode(2, 0)} & \texttt{1} \\
\hline
\end{tblr}
\end{table}

\subsection{Sample grader}

Сите повици до \texttt{encode} и \texttt{decode} ќе бидат во истото извршување на вашата програма. Исто така, \texttt{setup} ќе биде повиката само еднаш (наместо два пати, еднаш по извршување, како што е објаснато погоре.)

Влезот е само еден цел број -- $M$. Потоа ќе биде отпечатено вредноста на $N$ која \texttt{setup} ја вратила. Потоа ќе ги повика функциите \texttt{encode} и \texttt{decode} во овој редослед $T$ пати, со случајно генерирани броеви од $0$ до $N-1$ и случајно одбрани два броеви од трите кои ги вратила функцијата \texttt{encode} за да се пратат на \texttt{decode} (и во кој редослед). Ќе биде отпечатена порака која ви кажува дека вашето решение е погрешно.


\end{document}
